\begin{abstract}
\noindent Isharati is an AI-powered platform that brings the Qur'an and the Sunnah to Deaf and hard-of-hearing people in
their first language, sign language, and gives preachers and Islamic content creators integrated tools to reach them
without barriers. It answers a wide demographic and human need: more than 70 million Deaf people live worldwide (more
than 15 million of them Muslims), and they face a double linguistic and cultural isolation, because religious content
reaches them written or spoken in a language that is not theirs, while the signed content that exists is scarce,
scattered and without documented sources. Preachers, for their part, lack translation tools built for religious text:
tools that respect the precision of the scripture and avoid both unfounded interpretation and spelling Islamic terms
letter by letter instead of using their established signs.

To close this gap, Isharati brings education, answers and content creation together in one platform, built entirely
on documented sources. \emph{Verified question and answer} accepts a question in Arabic, English, Turkish or Urdu,
retrieves the relevant verses and hadith by retrieval-augmented generation (RAG), and accepts no sentence that its
source does not support; questions asking for a religious ruling are referred to scholars. The \emph{Academy} offers
interactive paths for learning the letters and religious terms in four sign languages, ending with a certificate of
completion, alongside a \emph{Qur'an and hadith corpus} of about 8,000 signed Qur'an segments (about 45 hours, from five
recognised sources such as the King Fahd Complex and the Al-Kharj curriculum) and more than 300 labelled hadith.

For institutions and preachers, the \emph{Content Studio} turns texts, audio or video lessons into signing at once,
with control over the signing speed, a choice of avatars in modest dress, and video export.

The sign lexicons cover most of the scripture: the share of content words in the Qur'an and hadith that the
lexicon of each language signs is 83.1\% for English to American Sign Language (3,677 signs), 75.7\% for Turkish to
Turkish Sign Language (4,827 signs), 72.9\% for Arabic to Arabic Sign Language (6,663 signs; 74.2\% with proper names
shown by fingerspelling), and 72.1\% for Urdu to Indian and Pakistani Sign Language (10,233 signs). Against the signs of
Deaf signers, the Arabic matcher chooses the right sign for 98.3\% of the words of religious sentences, with no wrong
sign.

Isharati thus offers a double value: direct and safe access for Deaf people to their religion at any time, and a
multiplied reach for preaching institutions at low cost and high efficiency. Its aim is that every Deaf Muslim receives
the Qur'an and the Sunnah as others do: understood, documented, and in their own language, so that Islamic content in
sign language becomes an integral part of preaching rather than an exception.
\end{abstract}

\begin{Arabic}
\noindent\textbf{الملخّص.}
تُعد «إشارتي» منصة رائدة مدعومة بالذكاء الاصطناعي تهدف إلى جعل القرآن الكريم والسنة النبوية في متناول الصم وضعاف
السمع بلغتهم الأم (لغة الإشارة)، مع منح الدعاة وصُنّاع المحتوى الإسلامي أدوات متكاملة تُمكّنهم من الوصول إلى هذه
الفئة دون حواجز. يأتي هذا الابتكار استجابةً لتحدٍ ديموغرافي وإنساني واسع؛ حيث يعيش أكثر من 70 مليون أصم حول العالم
(من بينهم أكثر من 15 مليون مسلم أصم)، يواجهون عزلة لغوية وثقافية مضاعفة نظراً لأن المحتوى الديني يصلهم مكتوباً أو
مسموعاً بلغة ليست لغتهم، فضلاً عن ندرة المحتوى الإشاري المتفرّق وافتقاره للمصادر الموثقة. وفي المقابل، يعاني الدعاة
من غياب أدوات الترجمة المتخصصة التي تحترم دقة النصوص الشرعية وتتجنب مخاطر الاجتهاد الخاطئ أو تهجئة المصطلحات حرفياً
بدل إشاراتها المعتمدة.

ولسد هذه الفجوة الهيكلية، تجمع منصة «إشارتي» في بيئة رقمية واحدة حلولاً مبتكرة للتعليم، والإجابة، وصناعة المحتوى
استناداً إلى مصادر موثقة بالكامل. تتضمن المنصة ميزة «سؤال وجواب موثق» التي تتيح للمستخدم الاستعلام بلغات متعددة
(العربية، الإنجليزية، التركية، الأردية) لتسترجع المنصة الآيات والأحاديث ذات الصلة عبر تقنية التوليد المعزز
بالاسترجاع (RAG) دون أي قبول لجملة لا يدعمها مصدرها، مع إحالة أسئلة الفتوى تلقائياً لأهل العلم. كما توفر «الأكاديمية»
مسارات تفاعلية لتعليم الحروف والمصطلحات الدينية بأربع لغات تنتهي بشهادة إتمام، إلى جانب «مدونة القرآن والحديث» التي
تضم نحو 8,000 مقطع قرآني (قرابة 45 ساعة من خمسة مصادر معتمدة كَمجمع الملك فهد ومنهج الخرج) وأكثر من 300 حديث نبوي
موسوم.

ولتمكين المؤسسات والدعاة، تقدم المنصة «استوديو صناعة المحتوى» الذي يتيح تحويل النصوص أو الدروس الصوتية والمرئية إلى
إشارة فورية مع التحكم في سرعة الأداء واختيار الشخصيات الافتراضية ذات الزي المحتشم وتصدير الفيديو بسلاسة.

وتغطي المعاجم الإشارية معظم النصوص الشرعية: إذ تبلغ نسبة الكلمات ذات المعنى في القرآن والحديث التي يملك معجم كل لغة
إشارةً لها \textenglish{83.1\%} للإنجليزية إلى لغة الإشارة الأمريكية (3,677 إشارة)، و\textenglish{75.7\%} للتركية إلى لغة الإشارة التركية (4,827
إشارة)، و\textenglish{72.9\%} للعربية إلى لغة الإشارة العربية (6,663 إشارة؛ و\textenglish{74.2\%} مع عرض أسماء الأعلام بالتهجئة الإشارية)، و\textenglish{72.1\%}
للأردية إلى لغة الإشارة الهندية والباكستانية (10,233 إشارة). ومقارنةً بإشارات مترجمين صمّ فعليين، يختار النظام
العربي الإشارة الصحيحة لـ\textenglish{98.3\%} من كلمات الجمل الدينية دون أي إشارة خاطئة.

وبذلك تقدم «إشارتي» قيمة مضاعفة تتمثل في وصول مباشر وآمن للصم إلى دينهم في أي وقت، ومضاعفة الأثر الدعوي للمؤسسات
بتكلفة منخفضة وكفاءة عالية، لتسعى المنصة في النهاية إلى أن يتلقى كل أصم مسلم القرآن والسنة كما يتلقاهما غيره:
مفهومان، موثقان، وببلغته، ليصبح المحتوى الإسلامي الإشاري جزءاً أصيلاً من العمل الدعوي لا استثناءً.
\end{Arabic}

\noindent\textbf{Keywords:} sign language, Arabic Sign Language, American Sign Language, retrieval-augmented generation,
Qur'an, hadith, accessibility, Deaf Muslims, 3D avatar.
